\documentclass[12pt,a4paper]{report}

\usepackage{graphicx}
\usepackage{url}
\usepackage{listings}
\usepackage{xcolor}
\usepackage[font=small,labelfont=bf]{caption}
\usepackage{fancybox}
\usepackage{amsmath}
\usepackage{amssymb}
\usepackage{setspace}
\usepackage{pdfpages}
\usepackage{hyperref}
\usepackage{fancyhdr}
\usepackage{times}
\usepackage[left=1in,right=1in,top=1in,bottom=0.8in]{geometry}
\usepackage{rotating}
\usepackage{array}
\usepackage{enumitem}
\usepackage{amsfonts}
\usepackage{float}

\hypersetup{
    colorlinks=true,
    linkcolor=black,
    urlcolor=blue,
    citecolor=black
}

\lstset{
    basicstyle=\ttfamily\small,
    breaklines=true,
    frame=single,
    numbers=left,
    numberstyle=\tiny,
    keywordstyle=\color{blue},
    commentstyle=\color{gray},
    stringstyle=\color{red},
    showstringspaces=false
}

\makeatletter
\def\@makechapterhead#1{%
  \vspace*{5\p@}%
  {\parindent \z@ \raggedright \normalfont
    \ifnum \c@secnumdepth >\m@ne
       \LARGE\bfseries \space \thechapter. \space
    \fi
    \interlinepenalty\@M
    \LARGE \bfseries #1\par\nobreak
    \vskip 10\p@
  }}
\makeatother

\newcommand{\optionalimage}[3]{
\begin{figure}[H]
\centering
\IfFileExists{#1}
  {\includegraphics[width=#2\textwidth]{#1}}
  {\fbox{\parbox[c][6cm][c]{#2\textwidth}{\centering Screenshot placeholder: #1}}}
\caption{#3}
\end{figure}
}

\begin{document}
\newpage
\pagestyle{plain}
\pagestyle{empty}
\pagestyle{fancy}
\renewcommand{\headrulewidth}{0pt}
\fancypagestyle{plain}{}

\pagenumbering{gobble}
\thisfancyput(-0.0in,-10.0in){%
\setlength{\unitlength}{1in}\framebox(6.7,10.2)}

\begin{titlepage}
\begin{center}

\textbf{SCTR's Pune Institute of Computer Technology}\\
\vspace{0.2cm}
Dhankawadi, Pune\\
A.Y. 2025--26

\vspace{1cm}

\large \textbf{CCL MINI ``SAAS'' PROJECT REPORT}

\vspace{0.6cm}

\normalsize ON

\vspace{0.8cm}

\textbf{\Large ``Contest Tracker: A Cloud-Based Competitive Programming Contest Reminder and Analytics Platform''}

\vspace{1cm}

\textbf{Submitted By}\\
\vspace{0.3cm}

Prathamesh Jadhav\\
Class: TE-09\\
Roll No: 33144

\vspace{1cm}

\textbf{Under the guidance of}\\
\vspace{0.2cm}
Mr. Sachin Pande

\vspace{2cm}

\IfFileExists{pict_logo.png}
  {\includegraphics[width=4cm]{pict_logo.png}}
  {\fbox{\parbox[c][3cm][c]{4cm}{\centering PICT Logo}}}

\vspace{2cm}

\textbf{DEPARTMENT OF INFORMATION TECHNOLOGY}\\
PUNE INSTITUTE OF COMPUTER TECHNOLOGY\\
SR. NO 27, PUNE-SATARA ROAD, DHANKAWADI\\
PUNE -- 411 043.\\

\vspace{0.5cm}

AY: 2025--2026

\end{center}
\end{titlepage}

\newpage
\pagestyle{plain}
\begin{center}
    \begin{LARGE}
        \section*{Abstract}
    \end{LARGE}
\end{center}
\vspace{0.5cm}

{\normalsize
\setlength{\baselineskip}{1.1\baselineskip}

\noindent
Contest Tracker is a cloud-enabled Software-as-a-Service web application designed for competitive programmers who need a single place to discover, filter, bookmark, and track programming contests. The platform aggregates contest data from supported coding platforms such as Codeforces, CodeChef, LeetCode, AtCoder, and Kaggle through the CLIST API, stores the normalized contest information in a PostgreSQL database, and presents it through a responsive Next.js interface.

\vspace{0.3cm}

\noindent
The system provides user authentication, personalized contest reminders, a calendar-based contest view, and Codeforces performance analytics. Registered users can save contests to their personal reminder list, and the backend reminder worker sends email notifications one hour before the selected contest starts. The profile module allows users to store their Codeforces handle and visualize rating history using chart components.

\vspace{0.3cm}

\noindent
The project demonstrates core Cloud Computing and SaaS concepts through a decoupled architecture consisting of a React/Next.js frontend, an Express.js REST API, Prisma ORM, PostgreSQL storage, scheduled background jobs, JWT-based authentication, and third-party API integration. It emphasizes scalability, automation, usability, and practical value for students and developers preparing for competitive programming events.

\vspace{0.6cm}

\noindent
\textbf{Keywords:} SaaS, Next.js, Express.js, PostgreSQL, Prisma, Competitive Programming, CLIST API, JWT, Email Reminder, Cloud Computing
}
\addcontentsline{toc}{section}{Abstract}

\newpage
\section*{Abbreviations}
\addcontentsline{toc}{section}{Abbreviations}

\begin{tabular}{ll}
SaaS & Software as a Service \\
API  & Application Programming Interface \\
UI   & User Interface \\
JWT  & JSON Web Token \\
ORM  & Object Relational Mapper \\
DB   & Database \\
SMTP & Simple Mail Transfer Protocol \\
CI/CD & Continuous Integration and Continuous Deployment \\
CRUD & Create, Read, Update, Delete \\
REST & Representational State Transfer \\
\end{tabular}

\newpage
\thispagestyle{empty}
\fancyhead{}
\renewcommand{\headrulewidth}{0pt}

\addcontentsline{toc}{section}{Contents}
\tableofcontents

\newpage
\pagestyle{fancy}
\pagenumbering{arabic}
\renewcommand{\headrulewidth}{0.5pt}
\fancyhead[RO]{\textit{Contest Tracker SaaS Platform}}
\fancyhead[LO]{\textit{CC Mini Project Report}}
\renewcommand{\footrulewidth}{0.5pt}
\fancyfoot[RO]{\textit{Dept. of Information Technology}}
\fancyfoot[LO]{\textit{PICT, Pune}}
\fancypagestyle{plain}{}

\chapter{\centering{Introduction}}
\begin{normalsize}

\section{Purpose}

The purpose of this project is to develop a SaaS-based contest tracking platform for competitive programmers. The application allows users to view upcoming and past coding contests, filter contests by platform, search contest names, add personal reminders, view contests in a calendar format, and analyze Codeforces rating progress from a single web dashboard.

\section{Problem Statement}

Competitive programming contests are hosted across several platforms, including Codeforces, CodeChef, LeetCode, AtCoder, and Kaggle. Students and developers often miss important contests because the schedules are scattered across different websites. Manually checking every platform is time-consuming and unreliable. Existing calendars may show contest information, but they often lack personalized reminders, user-specific tracking, and integrated performance analytics.

\vspace{0.3cm}

\noindent
Therefore, there is a need for a centralized, cloud-based platform that can aggregate contest schedules, allow users to save important contests, send timely reminders, and provide useful analytics for coding performance.

\section{Objectives}

\begin{itemize}
    \item To design and develop a responsive web application for tracking programming contests.
    \item To aggregate contest data from multiple competitive programming platforms using the CLIST API.
    \item To implement secure user registration and login using password hashing and JWT authentication.
    \item To provide personalized contest reminders through authenticated user accounts.
    \item To send automated email reminders before contests using scheduled backend jobs.
    \item To build a calendar interface for visualizing upcoming and past contests.
    \item To integrate Codeforces analytics for visualizing user rating history.
    \item To demonstrate a cloud-ready, decoupled SaaS architecture.
\end{itemize}

\section{Scope}

The scope of the project includes a full-stack contest tracking system with a Next.js frontend, Express.js backend, PostgreSQL database, Prisma ORM, authentication, contest synchronization, reminder management, and analytics visualization. The system supports multiple users, stores personalized reminder choices, and periodically updates contest data from an external provider.

\vspace{0.3cm}

\noindent
The current implementation focuses on contest discovery and reminders for selected platforms. Future versions can include payment plans, custom notification channels, team calendars, custom contest lists, mobile push notifications, and additional analytics for platforms such as LeetCode and CodeChef.

\end{normalsize}

\chapter{\centering{Motivation}}
\begin{normalsize}

Competitive programming is an important part of technical preparation for many engineering students. Regular participation helps improve problem-solving ability, data structures knowledge, algorithmic thinking, and interview readiness. However, contests are distributed across many platforms, each with different schedules and interfaces.

\vspace{0.3cm}

\noindent
The motivation behind Contest Tracker is to reduce this friction. A student should be able to open one dashboard and immediately know which contests are upcoming, which contests were missed, which contests have been bookmarked, and how their performance is improving over time.

\vspace{0.3cm}

\noindent
The project also provides hands-on exposure to practical cloud application development. It combines REST APIs, database modeling, authentication, scheduled workers, third-party API integration, charts, and modern frontend development. These components closely resemble production SaaS systems used in real-world software engineering.

\section{Need of the System}

\begin{itemize}
    \item Contest schedules are spread across many websites.
    \item Users need a personal reminder mechanism for important events.
    \item Calendar-based visualization makes planning easier.
    \item Performance analytics helps users understand their growth.
    \item A cloud-hosted centralized system improves accessibility from any device.
\end{itemize}

\end{normalsize}

\chapter{\centering{System Architecture and Design}}
\begin{normalsize}

\section{High-Level Architecture}

Contest Tracker follows a decoupled client-server architecture. The frontend is built using Next.js and React, while the backend is implemented using Express.js. Data persistence is handled using PostgreSQL through Prisma ORM. External contest data is fetched from the CLIST API, and Codeforces rating data is fetched from the Codeforces API. Email notifications are sent through SMTP using Nodemailer.

\begin{center}
\begin{tabular}{|p{3.2cm}|p{4.2cm}|p{6.2cm}|}
\hline
\textbf{Layer} & \textbf{Technology} & \textbf{Responsibility} \\
\hline
Frontend & Next.js, React, Tailwind CSS & User interface, contest listing, filters, calendar, profile dashboard \\
\hline
Backend API & Node.js, Express.js & REST endpoints, authentication, reminders, contest sync, analytics proxy \\
\hline
Database & PostgreSQL & Persistent storage for users, contests, and reminders \\
\hline
ORM & Prisma & Schema definition and database operations \\
\hline
Authentication & bcryptjs, JWT & Password hashing and token-based protected routes \\
\hline
External APIs & CLIST API, Codeforces API & Contest aggregation and rating history retrieval \\
\hline
Background Jobs & node-cron & Hourly contest sync and periodic reminder checks \\
\hline
Notification & Nodemailer, SMTP & Email reminders one hour before contest start time \\
\hline
\end{tabular}
\end{center}

\section{System Workflow}

\begin{enumerate}
    \item The backend periodically fetches contest data from the CLIST API for supported platforms.
    \item Contest records are normalized and stored in PostgreSQL using Prisma upsert operations.
    \item The user opens the Next.js frontend and views upcoming contests from the API.
    \item The user can search contests, filter by platform, or switch between upcoming, reminder, and history views.
    \item If the user wants personalized reminders, they register or log in.
    \item The backend verifies credentials, hashes passwords during registration, and issues JWT tokens after login.
    \item Authenticated users can toggle reminders for selected contests.
    \item A scheduled worker checks every five minutes for contests starting in about one hour and sends reminder emails.
    \item The user can add a Codeforces handle in the profile page and view rating history as a line chart.
\end{enumerate}

\section{Database Design}

The database consists of three main entities: \texttt{User}, \texttt{Contest}, and \texttt{Reminder}. A user can create many reminders, and each reminder links one user with one contest. The \texttt{Reminder} model uses a compound unique constraint on \texttt{userId} and \texttt{contestId} to prevent duplicate reminders.

\begin{center}
\begin{tabular}{|p{3cm}|p{9.5cm}|}
\hline
\textbf{Entity} & \textbf{Description} \\
\hline
User & Stores account details such as email, hashed password, name, and coding handles. \\
\hline
Contest & Stores contest platform, name, URL, start time, end time, and duration. \\
\hline
Reminder & Stores a user's saved contest reminder and links users with contests. \\
\hline
\end{tabular}
\end{center}

\section{API Design}

\begin{center}
\begin{tabular}{|p{4.2cm}|p{2.2cm}|p{7cm}|}
\hline
\textbf{Endpoint} & \textbf{Method} & \textbf{Purpose} \\
\hline
\texttt{/api/auth/register} & POST & Create a new user account. \\
\hline
\texttt{/api/auth/login} & POST & Authenticate user and return JWT token. \\
\hline
\texttt{/api/contests} & GET & Fetch upcoming contests. \\
\hline
\texttt{/api/contests/past} & GET & Fetch recent past contests. \\
\hline
\texttt{/api/platforms} & GET & Fetch distinct available platforms. \\
\hline
\texttt{/api/user/reminders} & GET & Fetch reminders for the authenticated user. \\
\hline
\texttt{/api/user/reminders/toggle} & POST & Add or remove a contest reminder. \\
\hline
\texttt{/api/user/profile} & POST & Update the user's Codeforces handle. \\
\hline
\texttt{/api/user/analytics/codeforces} & GET & Fetch Codeforces rating history. \\
\hline
\texttt{/api/contests/sync} & POST & Manually trigger contest synchronization. \\
\hline
\end{tabular}
\end{center}

\end{normalsize}

\chapter{\centering{Implementation Details}}
\begin{normalsize}

\section{Frontend Implementation}

The frontend is implemented with Next.js and React. The main dashboard displays contests in a responsive card grid. Users can switch between upcoming contests, personal contests, and contest history. Search and platform filters make it easier to find relevant contests quickly.

\vspace{0.3cm}

\noindent
The frontend includes the following major pages:

\begin{itemize}
    \item \textbf{Home Page:} Displays contest cards, filters, countdown timers, and reminder buttons.
    \item \textbf{Authentication Page:} Provides login and registration forms.
    \item \textbf{Calendar Page:} Uses \texttt{react-big-calendar} to show contests in month, week, and agenda views.
    \item \textbf{Profile Page:} Stores the user's Codeforces handle and displays rating history using Recharts.
\end{itemize}

\section{Authentication Implementation}

The backend provides register and login endpoints. During registration, the user's password is hashed using \texttt{bcryptjs}. During login, the entered password is compared with the hashed password stored in the database. After successful authentication, the server signs a JWT containing the user ID and email. Protected endpoints verify this token through middleware.

\begin{lstlisting}[language=JavaScript, caption={JWT authentication middleware}]
const authHeader = req.headers.authorization;
if (!authHeader || !authHeader.startsWith('Bearer ')) {
  return res.status(401).json({ error: 'Unauthorized: No token provided' });
}

const token = authHeader.split(' ')[1];
const decoded = jwt.verify(token, process.env.JWT_SECRET);
req.user = decoded;
next();
\end{lstlisting}

\section{Contest Synchronization}

Contest data is synchronized from the CLIST API. The backend filters supported resources and stores contest data in PostgreSQL. Prisma's \texttt{upsert} operation is used so that existing contests are updated and new contests are inserted without duplication.

\begin{lstlisting}[language=JavaScript, caption={Contest upsert logic}]
await prisma.contest.upsert({
  where: { url: item.href },
  update: {
    name: item.event,
    startTime: new Date(item.start),
    endTime: new Date(item.end),
    duration: item.duration,
    platform: item.resource
  },
  create: {
    platform: item.resource,
    name: item.event,
    url: item.href,
    startTime: new Date(item.start),
    endTime: new Date(item.end),
    duration: item.duration
  }
});
\end{lstlisting}

\section{Reminder System}

The reminder feature allows authenticated users to add or remove contests from their personal list. The backend checks whether a reminder already exists for the user and contest. If it exists, it is removed; otherwise, a new reminder record is created.

\vspace{0.3cm}

\noindent
A background worker runs every five minutes. It checks for contests starting between 60 and 65 minutes from the current time. For every matching reminder, an email is generated and sent to the user through Nodemailer and the configured SMTP service.

\section{Analytics Module}

The profile page allows users to store their Codeforces handle. The frontend requests rating history from the backend, and the backend retrieves it from the Codeforces API. The frontend converts the returned data into chart points and renders a line chart using Recharts. This helps users track rating progress across contests.

\section{Security Considerations}

\begin{itemize}
    \item Passwords are stored only after hashing with bcrypt.
    \item Protected routes require JWT tokens in the Authorization header.
    \item CORS is configured using the frontend URL from environment variables.
    \item Sensitive values such as database URL, JWT secret, CLIST credentials, and SMTP credentials are stored in environment variables.
    \item The reminder relation prevents duplicate reminders through a database-level unique constraint.
\end{itemize}

\end{normalsize}

\chapter{\centering{Tools and Technologies Used}}
\begin{normalsize}

\begin{center}
\begin{tabular}{|p{3.8cm}|p{4.2cm}|p{6cm}|}
\hline
\textbf{Category} & \textbf{Tool / Technology} & \textbf{Purpose} \\
\hline
Frontend Framework & Next.js, React & Building responsive client-side pages and components. \\
\hline
Styling & Tailwind CSS, Global CSS & Layout, colors, responsive UI, and dashboard styling. \\
\hline
Backend Framework & Node.js, Express.js & REST API development and server-side logic. \\
\hline
Database & PostgreSQL & Persistent relational data storage. \\
\hline
ORM & Prisma & Database schema modeling and query abstraction. \\
\hline
Authentication & bcryptjs, jsonwebtoken & Secure password hashing and JWT sessions. \\
\hline
Scheduler & node-cron & Contest sync and reminder background jobs. \\
\hline
Email Service & Nodemailer, SMTP & Sending reminder emails. \\
\hline
HTTP Client & Axios, Fetch API & Third-party and internal API requests. \\
\hline
Charts & Recharts & Codeforces performance visualization. \\
\hline
Calendar & react-big-calendar, date-fns & Calendar view and date formatting. \\
\hline
Icons & lucide-react & User interface icons. \\
\hline
\end{tabular}
\end{center}

\section{Next.js}

Next.js is used for building the frontend application. It provides a modern React-based development environment, file-based routing, optimized builds, and support for interactive client components.

\section{Express.js}

Express.js is used for creating REST APIs. It handles authentication, contest retrieval, reminder operations, profile updates, analytics requests, and manual contest synchronization.

\section{Prisma and PostgreSQL}

Prisma provides a type-safe and structured way to interact with the PostgreSQL database. The schema defines users, contests, and reminders, and Prisma Client performs database operations from the backend.

\section{CLIST and Codeforces APIs}

The CLIST API is used for fetching contest schedules from multiple platforms. The Codeforces API is used for retrieving user rating history, which is visualized in the profile dashboard.

\end{normalsize}

\chapter{\centering{Output and Results}}
\begin{normalsize}

\section{Output Screenshots}

The following figures represent the expected working screens of the Contest Tracker platform. If the screenshot files are placed in the same folder as this LaTeX file with the specified names, they will automatically appear in the report.

\optionalimage{home_dashboard.png}{0.9}{Home dashboard with upcoming contests, filters, and countdown timers}

\optionalimage{auth_page.png}{0.75}{User login and registration page}

\optionalimage{calendar_page.png}{0.9}{Contest calendar showing events across supported platforms}

\optionalimage{profile_analytics.png}{0.9}{Profile page with Codeforces rating analytics}

\section{Functional Results}

\begin{itemize}
    \item The system successfully fetches and stores contest data from supported competitive programming platforms.
    \item Upcoming and past contests are displayed through a clean web dashboard.
    \item Users can search contests and filter them by platform.
    \item Authenticated users can add or remove personal contest reminders.
    \item Email reminders are scheduled one hour before selected contests.
    \item Calendar view provides a visual planning interface for contests.
    \item Codeforces rating history is displayed as a performance chart.
    \item JWT authentication protects user-specific features.
\end{itemize}

\section{Cloud and SaaS Characteristics}

\begin{itemize}
    \item \textbf{Multi-user support:} Multiple users can register and maintain separate reminder lists.
    \item \textbf{Centralized service:} Contest data is stored and served through a backend API.
    \item \textbf{Automation:} Scheduled jobs update contest data and send reminder emails.
    \item \textbf{Scalability:} The client-server architecture can be deployed on cloud platforms with managed PostgreSQL.
    \item \textbf{Accessibility:} Users can access the dashboard from any modern browser.
\end{itemize}

\section{Deployment URLs}

\begin{itemize}
    \item \textbf{Frontend URL:} \texttt{Add deployed frontend URL here}
    \item \textbf{Backend API URL:} \texttt{Add deployed backend URL here}
    \item \textbf{GitHub Repository:} \texttt{Add GitHub repository URL here}
\end{itemize}

\end{normalsize}

\chapter{\centering{Testing}}
\begin{normalsize}

\section{Test Cases}

\begin{center}
\begin{tabular}{|p{1.2cm}|p{4cm}|p{4.5cm}|p{3.2cm}|}
\hline
\textbf{Sr. No.} & \textbf{Test Scenario} & \textbf{Expected Result} & \textbf{Status} \\
\hline
1 & User registration with valid details & Account is created and JWT token is returned & Pass \\
\hline
2 & Login with valid credentials & User is authenticated and redirected to dashboard & Pass \\
\hline
3 & Login with invalid credentials & Error message is displayed & Pass \\
\hline
4 & Fetch upcoming contests & Upcoming contests are displayed in ascending order & Pass \\
\hline
5 & Search contest by name & Matching contests are shown & Pass \\
\hline
6 & Filter contests by platform & Only contests from selected platform are shown & Pass \\
\hline
7 & Toggle contest reminder & Reminder is added or removed for authenticated user & Pass \\
\hline
8 & Access reminders without token & Unauthorized response is returned & Pass \\
\hline
9 & View calendar page & Contests are displayed as calendar events & Pass \\
\hline
10 & Save Codeforces handle & Handle is stored and analytics are fetched & Pass \\
\hline
\end{tabular}
\end{center}

\section{Limitations}

\begin{itemize}
    \item Contest data depends on the availability and correctness of the CLIST API.
    \item Email delivery depends on SMTP configuration and provider limits.
    \item The current analytics module focuses mainly on Codeforces rating history.
    \item Browser local storage is used to persist the frontend session token.
\end{itemize}

\end{normalsize}

\chapter{\centering{Conclusion}}
\begin{normalsize}
\setlength{\baselineskip}{1.3\baselineskip}

Contest Tracker successfully demonstrates the design and implementation of a practical SaaS-based platform for competitive programming contest management. The project solves a real problem faced by students and developers by centralizing contest schedules, enabling personalized reminders, and providing performance insights through analytics.

\vspace{0.3cm}

\noindent
The application combines a modern Next.js frontend with an Express.js backend, PostgreSQL database, Prisma ORM, JWT authentication, scheduled background jobs, third-party contest APIs, and email notifications. This architecture reflects real-world cloud application development practices and provides a scalable base for future enhancements.

\vspace{0.3cm}

\noindent
Future improvements can include mobile push notifications, Google Calendar integration, additional analytics for LeetCode and CodeChef, AI-based contest recommendations, team contest planning, and subscription-based SaaS plans. Overall, the project provides both technical learning value and practical utility for competitive programmers.

\end{normalsize}

\newpage
\chapter{\centering{References}}
\begin{normalsize}
\setlength{\baselineskip}{1.3\baselineskip}

\section{References}

\begin{enumerate}
    \item Next.js Documentation -- \url{https://nextjs.org/docs}
    \item React Documentation -- \url{https://react.dev}
    \item Express.js Documentation -- \url{https://expressjs.com}
    \item Prisma Documentation -- \url{https://www.prisma.io/docs}
    \item PostgreSQL Documentation -- \url{https://www.postgresql.org/docs}
    \item CLIST API Documentation -- \url{https://clist.by/api/v4/doc/}
    \item Codeforces API Documentation -- \url{https://codeforces.com/apiHelp}
    \item Nodemailer Documentation -- \url{https://nodemailer.com/about/}
    \item Recharts Documentation -- \url{https://recharts.org}
    \item React Big Calendar Documentation -- \url{https://github.com/jquense/react-big-calendar}
\end{enumerate}

\end{normalsize}

\end{document}
